\section{Aprendizaje por imitación vs. aprendizaje por refuerzo: comparación práctica}

Esta clase la imparte como invitado Ashish Kumar y se centra en las aplicaciones prácticas del RL en robótica, en particular en locomoción y manipulación. El ponente comienza comparando el aprendizaje por imitación y el aprendizaje por refuerzo «que funcionan en la práctica».

\begin{importantbox}{Rasgos prácticos de IL vs. RL}
\begin{itemize}
    \item \textbf{Aprendizaje por imitación}:
    \begin{itemize}
        \item Es eficiente en datos y conviene en escenarios donde se pueden obtener demostraciones de alta calidad
        \item Su rendimiento máximo está limitado por la capacidad del demostrador
        \item Difícilmente supera el nivel humano
    \end{itemize}
    \item \textbf{Aprendizaje por refuerzo}:
    \begin{itemize}
        \item Puede descubrir políticas que superan a las humanas
        \item Suele requerir un entorno de simulación (sim-to-real transfer)
        \item El diseño de la recompensa es el principal desafío
        \item El entrenamiento es largo, pero la inferencia es eficiente
    \end{itemize}
\end{itemize}
\end{importantbox}

\subsection{Cuándo elegir RL}

\begin{knowledgebox}{Condiciones en que conviene el RL}
El RL tiene ventaja sobre el IL en los siguientes escenarios:
\begin{itemize}
    \item Se dispone de un \textbf{simulador} de alta calidad
    \item El objetivo de la tarea puede cuantificarse con una función de recompensa
    \item Se necesita \textbf{superar} el nivel de los demostradores humanos
    \item Los datos de demostración son difíciles de obtener (por ejemplo, en operaciones peligrosas)
    \item Se requiere \textbf{robustez} ante cambios en el entorno
\end{itemize}
\end{knowledgebox}

\subsection{Resumen de la sección}

IL y RL tienen cada uno sus escenarios de aplicación, y en los sistemas reales suelen combinarse (por ejemplo, IL warm-start + RL fine-tuning).

